\chapter{快速入门}
\label{ch:3}

本章内容：在单台开发机上用约 10 分钟跑通 EdgeFlow 完整链路：构建制品、启动云端与边缘、注册节点、下发 Pod、查看容器、查看设备数据、下发设备指令，并可选体验模型发布。适合第一次接触 EdgeFlow 的操作人员。

\section{准备工作}

开始前请确认：

\begin{itemize}
  \item 本机为 macOS 或 Linux；
  \item Docker 已安装且 daemon 运行中（Edged 的容器运行时）；
  \item 已安装 curl；源码构建还需 Go 1.26+ 与 Make；
  \item 已进入 EdgeFlow 仓库根目录（\texttt{cd edgeflow}）；
  \item 可选：mosquitto（MQTT 数据面），缺失时跳过不影响主链路。
\end{itemize}

\begin{lstlisting}
docker info        # 确认 Docker daemon 运行
go version         # 确认 Go 版本（源码构建时需要）
\end{lstlisting}

\section{第一步：构建制品}

\begin{lstlisting}
make build
\end{lstlisting}

构建产物位于 \texttt{bin/}：\texttt{bin/cloudcore}、\texttt{bin/edgecore}。也可以直接使用发布目录 \texttt{release/v0.7.0/} 中的预编译二进制（获取方式见第 2 章 \autoref{ch:2}）。

\section{第二步：启动云端 cloudcore}

打开终端 A，执行：

\begin{lstlisting}
./bin/cloudcore
\end{lstlisting}

预期日志输出 \texttt{HTTP server listening on :8080}。默认监听：REST API 端口 8080、CloudHub WebSocket 端口 10000。端口可通过环境变量覆盖（\texttt{EDGEFLOW\_CLOUDCORE\_PORT}、\texttt{EDGEFLOW\_CLOUDCORE\_HUB\_PORT}），命令行参数 \texttt{--port} 优先级更高。

验证健康检查：

\begin{lstlisting}
curl http://127.0.0.1:8080/healthz
\end{lstlisting}

预期返回 HTTP 200 与 \texttt{\{"status":"ok","version":\{...\}\}}。

\section{第三步：启动边缘 edgecore}

打开终端 B，执行：

\begin{lstlisting}
EDGEFLOW_EDGECORE_NODE_ID=node-1 \
EDGEFLOW_EDGECORE_CLOUD_ADDR=ws://127.0.0.1:10000/v1/edge \
./bin/edgecore
\end{lstlisting}

说明：

\begin{itemize}
  \item \texttt{EDGEFLOW\_EDGECORE\_NODE\_ID}：节点 ID。不设置时默认为 \texttt{edge-<主机名>}，仅允许字母、数字、点、连字符与下划线；
  \item \texttt{EDGEFLOW\_EDGECORE\_CLOUD\_ADDR}：云端 CloudHub 地址（明文通道用 \texttt{ws://}；启用 mTLS 时用 \texttt{wss://}，见第 2 章 \autoref{ch:2}）；
  \item edgecore 启动后自动完成节点注册、云边连接，并随进程启动内置模拟温湿度传感器 \texttt{sensor-01}（mock\_sensor），无需额外配置。
\end{itemize}

\section{第四步：验证节点注册}

回到终端 A，查询节点列表：

\begin{lstlisting}
curl http://127.0.0.1:8080/api/v1/nodes
\end{lstlisting}

预期返回 JSON 数组，包含 \texttt{nodeID} 为 \texttt{node-1}、\texttt{status} 为 \texttt{Ready} 的节点。若长时间未出现，请检查终端 B 的 edgecore 日志与网络连通性。

\section{第五步：下发第一个 Pod（nginx）}

\begin{lstlisting}
curl -X POST http://127.0.0.1:8080/api/v1/nodes/node-1/podsync \
  -H 'Content-Type: application/json' \
  -d '{"operation":"add","pod":{"name":"nginx","namespace":"default","image":"nginx:1.25-alpine","replicas":1}}'
\end{lstlisting}

预期返回 \texttt{\{"status":"ok","acked":true\}}，表示边缘节点已确认。镜像首次拉取可能需要一些时间。

如需限制容器资源（可选，v0.2.0 起支持），在 \texttt{pod} 中追加 \texttt{resources} 字段（K8s 风格，零值表示不限制）：

\begin{lstlisting}
curl -X POST http://127.0.0.1:8080/api/v1/nodes/node-1/podsync \
  -H 'Content-Type: application/json' \
  -d '{"operation":"add","pod":{"name":"nginx","namespace":"default","image":"nginx:1.25-alpine","replicas":1,"resources":{"cpuRequest":"100m","cpuLimit":"250m","memoryRequest":"64Mi","memoryLimit":"128Mi"}}}'
\end{lstlisting}

资源语义（v0.2.0）：请求量大于上限（\texttt{request} > \texttt{limit}）返回 400；超出节点超卖容量返回 409（\texttt{EDGEFLOW\_RESOURCE\_EXHAUSTED}，不落盘不建容器）；格式非法（如 NaN/Inf/前导 \texttt{+}）同样返回 400。

\section{第六步：查看容器与 Pod 状态}

在边缘侧查看容器（本机即边缘）：

\begin{lstlisting}
docker ps --filter label=edgeflow.pod
\end{lstlisting}

预期看到容器 \texttt{edgeflow-default-nginx-0} 处于运行状态。

查询云端 Pod 状态：

\begin{lstlisting}
curl http://127.0.0.1:8080/api/v1/pods
\end{lstlisting}

预期返回 \texttt{phase} 为 \texttt{Running} 的 Pod 记录。

\begin{note}资源漂移说明（v0.2.0 起）：在边缘侧直接修改容器资源限制（如 \texttt{docker update}）会被 Edged 判定为与期望状态不一致并自动重建容器（每轮最多重建 1 个），资源以云端下发为准。\end{note}

\section{第七步：查看设备数据}

\begin{lstlisting}
curl http://127.0.0.1:8080/api/v1/devices
\end{lstlisting}

预期返回设备 \texttt{sensor-01}，\texttt{properties} 中包含 \texttt{temperature} 与 \texttt{humidity} 两个属性（周期上报的实时值）。设备默认上报周期为 30 秒，可设置 \texttt{EDGEFLOW\_EDGECORE\_DEVICE\_REPORT\_INTERVAL=3s} 加速观察。

\section{第八步：下发设备指令}

向模拟传感器下发目标温度指令：

\begin{lstlisting}
curl -X POST http://127.0.0.1:8080/api/v1/nodes/node-1/device-command \
  -H 'Content-Type: application/json' \
  -d '{"deviceName":"sensor-01","property":"targetTemp","value":25}'
\end{lstlisting}

预期返回 \texttt{\{"status":"ok","acked":true\}}。再次查询设备：

\begin{lstlisting}
curl http://127.0.0.1:8080/api/v1/devices
\end{lstlisting}

预期 \texttt{sensor-01} 的 \texttt{desired} 中写入 \texttt{\{"targetTemp":25\}}，表明指令已被边缘执行并写回云端设备状态。

\section{体验模型发布（可选，v0.7.0）}

以下步骤在已完成第八步的同一环境上体验模型仓库与灰度发布。该能力为云端内置（v0.7.0），边缘零代码改动；发布器复用既有 podsync 与 config-sync 经可靠投递下发，旧版 edgecore 直接可用。前置条件：\texttt{node-1} 处于 \texttt{Ready} 状态（百分比发布的分母为创建时刻 Ready 节点数，0 台会返回 422）。

\subsection{第一步：注册模型}

\begin{lstlisting}
curl -X POST http://127.0.0.1:8080/api/v1/models \
  -H 'Content-Type: application/json' \
  -d '{"name":"defect-detector","description":"缺陷检测模型","type":"detection"}'
\end{lstlisting}

预期返回 200 与完整 Model 对象（含 \texttt{createdAt}/\texttt{updatedAt}）。模型名全局唯一（重复创建返回 409）。

\subsection{第二步：登记版本}

\begin{lstlisting}
curl -X POST http://127.0.0.1:8080/api/v1/models/defect-detector/versions \
  -H 'Content-Type: application/json' \
  -d '{"version":"v1.2.0","mirror":"registry.example.com/edgeflow/models/defect-detector:v1.2.0","sha256":"sha256:9f86d081884c7d659a2feaa0c55ad015a3bf4f1b2b0b822cd15d6c15b0f00a08","archs":["amd64","arm64"],"metadata":{"threshold":"0.8"}}'
\end{lstlisting}

预期返回 200 与 ModelVersion（\texttt{status} 为 \texttt{draft}）。\textquotedblleft Tag 即版本\textquotedblright ：\texttt{mirror} 必须带 tag，\texttt{sha256} 为镜像防篡改登记（\texttt{\textasciicircum{}sha256:[0-9a-f]\{64\}\$}）。

\subsection{第三步：激活版本}

\begin{lstlisting}
curl -X POST http://127.0.0.1:8080/api/v1/models/defect-detector/versions/v1.2.0/activate
\end{lstlisting}

预期返回 200；版本状态 \texttt{draft} → \texttt{active}（自动降级旧 active 版本）。发布目标必须是 active 版本。

\subsection{第四步：全量发布}

\begin{lstlisting}
curl -X POST http://127.0.0.1:8080/api/v1/models/defect-detector/releases \
  -H 'Content-Type: application/json' \
  -d '{"version":"v1.2.0","target":{"type":"percentage","percentage":100}}'
\end{lstlisting}

预期返回 202 与 release 对象（\texttt{status} 为 \texttt{pending}，\texttt{targetNodes} 已物化）。发布为异步受理：可用 \texttt{GET /api/v1/models/defect-detector/releases/<releaseID>} 跟踪执行进度，全部目标节点 \texttt{deployed} 后进入 \texttt{succeeded}。

\subsection{第五步：查询部署影子}

\begin{lstlisting}
curl http://127.0.0.1:8080/api/v1/models/defect-detector/deployments
\end{lstlisting}

预期返回每台目标节点的\textquotedblleft 版本—节点—时间\textquotedblright 记录（\texttt{model}/\texttt{version}/\texttt{mirror}/\texttt{releaseID}/\texttt{updatedAt}），即云端写穿的部署影子。

说明：模型发布为云端能力——发布器自动执行 podsync（Pod 名 \texttt{edgeflow-model-<模型名>}、命名空间固定 \texttt{edgeflow}、replicas=1）与 config-sync（ConfigMap 携带 model/version/mirror/sha256/type/releasedAt 与版本参数），两步均被边缘确认后写穿部署影子。注意：\textbf{发布成功 ≠ 镜像可用}——拉取发生在边缘，请以试点节点 PodStatus（Running/CrashLoop）为准；灰度节奏建议先白名单 1 台试点 → 小批（\texttt{batchSize}+\texttt{pauseBetween}）→ 全量；生产环境建议 embed/外部 etcd 模式（纯内存模式下发布任务与部署影子重启丢失）。

\section{一键 Demo（可选）}

仓库提供端到端演示脚本，自动完成上述全部步骤并输出 \texttt{DEMO PASS}：

\begin{lstlisting}
bash examples/demo.sh
\end{lstlisting}

脚本自动完成：构建、挑选空闲端口、启动 cloudcore 与 edgecore（运行时数据落在临时目录）、验证节点注册、下发 nginx Pod、验证容器与 Pod 状态、验证设备数据流、下发设备指令、可选 MQTT 数据面验证，最后自动清理资源。脚本幂等，可重复运行；设置 \texttt{EDGEFLOW\_DEMO\_SKIP\_BUILD=1} 可跳过构建步骤。

\section{清理}

\begin{lstlisting}
# 1. 回收 Pod（边缘 Edged 会自动删除容器）
curl -X POST http://127.0.0.1:8080/api/v1/nodes/node-1/podsync \
  -H 'Content-Type: application/json' \
  -d '{"operation":"delete","pod":{"name":"nginx","namespace":"default"}}'

# 2. 停止 cloudcore 与 edgecore（终端 A/B 按 Ctrl+C，优雅退出）

# 3. 兜底清理容器与本地元数据（可选）
docker rm -f $(docker ps -aq --filter label=edgeflow.pod) 2>/dev/null
rm -f data/edgeflow.db
\end{lstlisting}

\section{常见问题}

\begin{tabularx}{\textwidth}{|l|X|}
  \hline
  \textbf{现象} & \textbf{处理} \\
  \hline
  节点未注册或状态非 Ready & 查看 edgecore 日志；确认 CloudHub 端口（默认 10000）可达、云边地址配置正确 \\
  \hline
  podsync 返回 504 或长时间未确认 & 边缘未确认下发：确认 edgecore 正在运行，且下发路径中的节点 ID 与注册的 nodeID 完全一致 \\
  \hline
  容器未创建 & 确认 Docker daemon 运行（\texttt{docker info}）；镜像首次拉取需要时间 \\
  \hline
  设备无数据上报 & 上报周期默认 30 秒，可设置 \texttt{EDGEFLOW\_EDGECORE\_DEVICE\_REPORT\_INTERVAL=3s} 加速观察 \\
  \hline
\end{tabularx}
